\documentclass[12pt]{article}
\usepackage[utf8]{inputenc}
\usepackage[T1]{fontenc}
\usepackage[margin=1in]{geometry}
\usepackage{setspace}
\usepackage{xcolor}
\usepackage[colorlinks=true, urlcolor=blue!60!black, linkcolor=blue!60!black]{hyperref}
\setlength{\parindent}{0em}
\setlength{\parskip}{0.45em}
\renewcommand{\baselinestretch}{1.0}
\begin{document}
\begin{center}
{\Large \textit{Citizens United v. Federal Election Commission}, 558 U.S. 310 (2010)\par}
\vspace{0.3em}
{\normalsize Case Brief --- One-Page Edition\par}
\end{center}

\textbf{Memory Jogger.} No bans on corporate or union independent political expenditures; disclaimers and disclosure are fine (\href{https://www.courtlistener.com/opinion/1741/citizens-united-v-federal-election-commission/}{CourtListener}).

\textbf{Facts.} Citizens United, a nonprofit corporation, made \textit{Hillary: The Movie}, attacking Senator Hillary Clinton during the 2008 primaries, and planned video-on-demand distribution and TV ads. McCain--Feingold barred corporate treasury spending on ``electioneering communications'' near elections, leaving only PACs.

\textbf{Procedural History.} A three-judge district court ruled for the FEC; the Supreme Court ordered reargument on overruling \textit{Austin}.

\textbf{Votes.} 5--4 on the expenditure ban (Kennedy, with Roberts, Scalia, Alito, Thomas); 8--1 on disclosure (Thomas dissenting). Stevens, with Ginsburg, Breyer, and Sotomayor, dissented in part.

\textbf{Holding.} The government may not suppress political speech on the basis of the speaker's corporate identity. Independent expenditures do not create \textit{quid pro quo} corruption, so the ban fails strict scrutiny; \textit{Austin} and part of \textit{McConnell} are overruled. Disclosure survives exacting scrutiny.

\textbf{Reasoning.} Section 441b was an outright ban, and the PAC option was no cure. Government has no interest in equalizing voices, and ingratiation and access are not corruption. Roberts: ``There is a difference between judicial restraint and judicial abdication.'' Scalia: the Amendment protects speech, not categories of speakers. Stevens: corporations cannot vote or run for office, and the Framers saw corruption more broadly. Thomas would also have struck disclosure.

\textbf{Examples.} \emph{Same-Side:} a state corporate-expenditure ban falls. \emph{Opposite-Side:} the ban on direct corporate contributions survives. \emph{Fence-Sitter:} spending by foreign-owned U.S. subsidiaries.

\textbf{Critique.} Critics say Super PACs and ``dark money'' disprove the no-corruption premise; defenders say the law silenced nonprofits and small businesses.

\textbf{Key Quotation.} ``For the reasons explained above, we now conclude that independent expenditures, including those made by corporations, do not give rise to corruption or the appearance of corruption.'' --- Justice Kennedy, 558 U.S. at 357.

\end{document}
